\section{Introduction}
\label{sec:intro}

Short-horizon complication prediction from irregular electronic health record
(EHR) streams has been driven by two lines of work with an unresolved
intersection. On one side, sequence models---GRU-D
\citep{che2018recurrent}, STraTS \citep{tipirneni2022strats}, visit-level
Transformers such as BEHRT \citep{li2020behrt}, TransformEHR
\citep{yang2023transformehr}---operate on raw measurement streams and
extract clinically meaningful signal without explicit medical knowledge,
paying the cost that their attributions resolve to raw variables rather
than to clinical concepts. On the other, symbolic pipelines built on
knowledge-based temporal abstraction (KBTA) \citep{shahar1997framework}
convert raw observations into named clinical intervals---states, trends,
events, contexts---that support downstream mining (e.g., KarmaLego
\citep{moskovitch2015karmalego}) or scoring, at the cost of a curated
knowledge base and typically at some loss of predictive power against
deep-learning alternatives. The interval-based Transformer INTERVenE-Enc
\citep{oded2027intervene} sits between the two, taking a KBTA interval
stream as its token sequence and outperforming both symbolic pipelines
and raw-stream neural baselines on a MIMIC-IV diabetic cohort while
supporting concept-level explanations through gradient-decoupled
sparse-transcoder attribution.

Two questions remain unresolved after that MIMIC-IV study. First, whether
its central architectural claim---that a Transformer over KBTA intervals
outperforms Transformers over the raw stream---transfers to a second cohort
with a different underlying data model and a different clinical knowledge
base. Ablations on MIMIC-IV \citep{oded2027intervene} suggested that
replacing KBTA intervals with knowledge-free $\sigma$-bin discretizations
left the AUPRC essentially unchanged: on that data the knowledge base
contributed interpretability but not accuracy. Whether this generalizes to
a network the knowledge base was actually authored for is the question the
present study asks.

Second, whether a measured signal of \emph{care quality} carries predictive
information beyond the record it is computed from. In parallel work
\citep{oded2027kineret-jbi}, the same abstraction engine was used to score
adherence to fourteen recommendations from a national inpatient
hyperglycemia guideline \citep{hochberg2023inpatient} across a
six-hospital cohort, yielding continuous per-recommendation compliance
scores at every point along an admission. The clinical study found
above-median adherence associated with lower odds of $30$-day mortality
after propensity weighting (OR $=0.557$, 95\% CI $[0.514, 0.603]$) and
identified the recommendations most strongly linked to outcomes as those
with the largest implementation gaps. Whether feeding those scores into a
predictive model changes what the model does---beyond the prediction the
underlying abstractions already support---was left open.

This paper answers both questions on the Kineret cohort. Because we control
the input representation while everything else is held constant, the same
comparison also settles a subsidiary point: whether the marginal
contribution of clinical knowledge to predictive accuracy on data the
knowledge base was authored for is the same as, or different from, its
marginal contribution on data it was not.

\subsection*{Contributions}

\begin{itemize}\setlength{\itemsep}{.15em}
\item A \emph{seven-arm predictive ladder}
      (Section~\ref{sec:methods-ladder}) that holds cohort, labels, split
      and evaluation window constant and varies only the input
      representation, so the effect of the architecture, the effect of the
      knowledge base, and the effect of the quality-assessment signal each
      appear as a single ordered contrast on the same numbers.
\item A demonstration on $4{,}686$ held-out test admissions that
      INTERVenE-Enc on KBTA intervals reaches a support-weighted AUPRC of
      $0.527$ against $0.472$ for the strongest raw-stream baseline
      (ss-STraTS), with non-overlapping $95\%$ confidence intervals, and
      that replacing KBTA with $\sigma$-bin discretization costs $0.073$
      AUPRC and drops the arm to the logistic-regression floor---a lift
      that was within noise on MIMIC-IV.
\item A first empirical answer to whether measured guideline adherence
      helps prediction: the QA signal is picked up by the model
      (sparse-transcoder attribution places compliance-pattern tokens
      among the top drivers for seven of eight outcomes) but does not
      move aggregate metrics, because it is a deterministic derivation of
      the same abstraction stream the model is already reading.
\item An \emph{event-reconciliation pipeline} that harmonizes the two
      source tables' disagreeing complication tallies into one canonical
      event set, so every arm is scored on identical labels and identical
      patients---a wiring check that fails loudly if not.
\end{itemize}

The remainder of the paper is organized as follows.
Section~\ref{sec:background} covers the necessary background on KBTA,
interval-based Transformers, and inpatient-diabetes guidelines; a full
treatment of INTERVenE-Enc's architecture is in the companion methodological
paper \citep{oded2027intervene} and a full treatment of the guideline
adherence study is in \citep{oded2027kineret-jbi}. Section~\ref{sec:methods}
describes the cohort, the seven arms and the evaluation protocol.
Section~\ref{sec:results} reports the target support, the ladder, the QA
delta, per-outcome results, and the time-head errors. Section~\ref{sec:discussion}
discusses transferability, the knowledge-base lift here versus on MIMIC-IV,
and why the QA signal does not shift aggregates. Section~\ref{sec:conclusion}
concludes. Reproducibility artefacts are listed at the end of the paper.
